% Imported from topics/video-generation/seedance-2.5/adapter-contract.md
\chapter{新增模块的维度与接入契约}\label{seedance-m21}
\moduleinfo{M21}{候选设计与工程选型；组合效果待验证}{\hyperref[src-ltx-2]{\texttt{ltx-2}}、\hyperref[src-ltx-original-code]{\texttt{ltx-original-code}}、\hyperref[src-ltx-original-model]{\texttt{ltx-original-model}}、\hyperref[src-vace]{\texttt{vace}}}

前面分别讲了注意力和音视频交换，这里固定一个公开替代基座，逐层核对候选条件支路的维度与插入顺序。

已有音视频主干之后，怎样增加控制而不把原来所有接口重做？先把原权重的输入约定固定，再定义新条件从哪里读、以什么维度加回来。下面以 LTX-2 原系列 19B 开发模型逐层写出这种关系。原生配置来自公开元数据，新增宽度、层位与训练方式则是未验证的候选。

\section{先固定一套确切的主干}\label{seedance-m21-ux5148ux56faux5b9aux4e00ux5957ux786eux5207ux7684ux4e3bux5e72}

选用 \codepath{ltx-2-19b-dev.safetensors}，保留配套视频/音频codec、Gemma文本特征、连接器、归一化和采样器。模型配置为48个联合块，视频宽度4096、32头且每头128维，音频宽度2048、32头且每头64维。两流输入与输出token均为128维。

这组权重使用split RoPE、双精度频率计算，时间调制倍率1000，AV交互调制倍率1000。构造器提供的默认值不能代替权重配置；新参考支路也不改写原有RoPE。

{\def\LTcaptype{table} % do not increment counter
\small\setlength{\tabcolsep}{3pt}\renewcommand{\arraystretch}{1.18}
\begin{longtable}[]{@{}
  >{\raggedright\arraybackslash}p{(\linewidth - 4\tabcolsep) * \real{0.1900}}
  >{\raggedright\arraybackslash}p{(\linewidth - 4\tabcolsep) * \real{0.3900}}
  >{\raggedright\arraybackslash}p{(\linewidth - 4\tabcolsep) * \real{0.4200}}@{}}
\toprule\noalign{}
\begin{minipage}[b]{\linewidth}\raggedright
对象
\end{minipage} & \begin{minipage}[b]{\linewidth}\raggedright
进入模型前
\end{minipage} & \begin{minipage}[b]{\linewidth}\raggedright
进入联合块后
\end{minipage} \\
\midrule\noalign{}
\endhead
\bottomrule\noalign{}
\endlastfoot
视频目标 & \(B\times N_v\times128\) & \(B\times N_v\times4096\) \\
音频目标 & \(B\times N_a\times128\) & \(B\times N_a\times2048\) \\
文本连接器输出 & \(B\times L\times3840\) & 原生投影分别到4096/2048 \\
视频坐标边界 & \(B\times3\times N_v\times2\) & 原生时空位置编码 \\
音频坐标边界 & \(B\times1\times N_a\times2\) & 原生时间位置编码 \\
每token噪声状态 & \(B\times N_m\) & 原生AdaLN与AV调制 \\
\end{longtable}
}

\(B\)为batch，\(N_v,N_a,L\)为视频、音频、文本长度。坐标最后一维保存原生位置网格的起止；视频时间轴按帧率换算为秒，空间轴保留原生尺度，不能把latent索引直接当成秒。这些位置边界也不等于codec的完整感受野，编辑影响范围需要另行扩张。

\section{新增模块放在什么位置}\label{seedance-m21-ux65b0ux589eux6a21ux5757ux653eux5728ux4ec0ux4e48ux4f4dux7f6e}

原联合块先作流内 self-attention 与文本读取，再进行双向 AV 交换，最后各自过 FFN；交换两边都读更新前状态。候选新模块接在第 4、8、12，直到第 48 个原块之后，共 12 处。这样原块和新支路的先后关系就确定了。

每处再按下面的顺序计算，不能在不同版本中只保留相同模块名称，却悄悄交换执行位置：

\begin{verbatim}
原生AV块
  → 对齐源条件/结构条件残差，更新视频流
  → 视觉参考读取，更新视频流
  → 声音参考读取，更新音频流
  → 下一个原生AV块
\end{verbatim}

加进来的条件会在下一原块的 AV 交换里继续传播。最后一处则已经在第 48 块后，残差直接走向输出头，不应再凭空假设还有一次原生 AV 交换。这个差别从调用顺序就能看出。

\section{参考支路的维度与共享规则}\label{seedance-m21-ux53c2ux8003ux652fux8defux7684ux7ef4ux5ea6ux4e0eux5171ux4eabux89c4ux5219}

为使读取支路有共同维度，摘要记忆暂取 512 维，视觉与声音分别池化。每图 32、动作参考 64、音色参考 32 个摘要位置，只是起始预算；精确时序音轨仍保留完整条件，不能用摘要位置替代。

每个插入位置独立拥有Q/K/V与输出投影，8个头、每头64维：

\[
Q_v=\operatorname{RMSNorm}(h_v)W_q^v,\quad
K_r=RW_k^v,\quad V_r=RW_v^v,
\]

\[
h'_v=h_v+\tanh(g_v)
\operatorname{Attn}(Q_v,K_r,V_r;B)W_o^v.
\]

这里可以再检查一次起始梯度。若残差记为 \(\tanh(g)f_\phi\)，则
\[
\partial_g[\tanh(g)f_\phi]_{g=0}=f_\phi,
\qquad
\partial_\phi[\tanh(g)f_\phi]_{g=0}=0.
\]
因此正常初始化的支路先让门控学习，门控离开零以后才进一步更新支路参数。若 \(f_\phi\) 也被置成零，第一项便可能一起消失，不能只凭“零初始化稳定”就把所有因子全置零。

\(W_q^v\)为 \(4096\times512\)，\(W_o^v\)为 \(512\times4096\)，\(W_k^v,W_v^v\)为 \(512\times512\)。音频支路将4096替换成2048。素材编码和池化器跨12处共享，读取投影与门控逐处独立；参考K/V因此可以按层缓存，不能误认为所有层拥有同一K/V。

为什么只把门控置零？门控为零时初始残差严格为零，正常初始化的投影先提供一条非零候选路径，门控因而能获得梯度；门控更新后梯度再进入支路。若输出投影也同时全零，两个因子就可能互相挡住起始更新。因此输出线性层不加 bias，其余投影正常初始化。

示意实现见 \hyperref[code-gated-condition]{gated\_condition.py}。它省略多头布局、时间调制与区域kernel，仅说明零门控和维度关系。

\section{源视频支路保留空间网格}\label{seedance-m21-ux6e90ux89c6ux9891ux652fux8defux4fddux7559ux7a7aux95f4ux7f51ux683c}

设编辑mask \(M=1\) 为允许变化，源视频拆为 \(F_c=F\odot M\) 与 \(F_k=F\odot(1-M)\)。两路独立经过原视频VAE，并将mask按感受野对齐至目标网格：

\[
S_0=P_{257\rightarrow512}
[E_v(F_c);E_v(F_k);D(M)].
\]

128＋128＋1得到257通道；这里的一通道mask是候选简化，不声称复制VACE的全部mask重排。\(S_0\)保留 \(T_v,H_v,W_v\)，与每个目标视频token逐点对应。

为避免复制完整19B主干，候选源条件塔采用512宽的轻量残差块：逐点通道混合、时间深度可分离卷积与空间深度可分离卷积。联合生成版可读条件邻域；因果续写版必须使用左侧时间卷积和同样的块边界。每处由独立的 \(512\rightarrow4096\) 投影与零门控把条件特征加回视频流。

结构条件也先注明自身类型。depth、normal、pose 与相机射线分别投影到 512 维，附有效位和置信度，再同源条件塔融合。深度的物理数值不能按普通 RGB 归一化，否则输入表面形状相同，含义已经改变。

源塔在首版只读干净条件与固定任务标签，不读当前带噪目标，所以这些条件残差在推理步之间可以固定。若条件运动本来就与目标位置不对齐，则要改用结构对齐或参考记忆，不能为了逐点相加硬造一个错误对应。

\section{三个状态开关必须分开}\label{seedance-m21-ux4e09ux4e2aux72b6ux6001ux5f00ux5173ux5fc5ux987bux5206ux5f00}

\begin{enumerate}
\def\labelenumi{\arabic{enumi}.}
\tightlist
\item
  \codepath{present}：这个模态或条件是否存在。
\item
  \codepath{process_hidden}：该流在Transformer中是否继续更新隐藏状态。
\item
  \codepath{denoise_update}：采样器是否更新它的latent。
\end{enumerate}

保留音轨时，音频latent保持干净、噪声时间为0，但音频流仍可处理隐藏状态并参与AV交换；采样器不更新其latent。仅将原生 \codepath{enabled=False} 不能表达这一完整策略，因为它还会跳过该流的自注意力、文本与FFN路径。

音频输入 latent 固定后，每步仍从它重新构造隐藏状态。因为深层声音特征可能已经读到当步视频，视频变化时它也变化。因此 codec 结果能缓存，并不能推出整条声音 Transformer 都能缓存。

\section{自定义策略需要改到哪些工程接口}\label{seedance-m21-ux81eaux5b9aux4e49ux7b56ux7565ux9700ux8981ux6539ux5230ux54eaux4e9bux5de5ux7a0bux63a5ux53e3}

已有 \codepath{TextToVideoStrategy} 可作为共同sigma音视频训练基线；\codepath{VideoToVideoStrategy} 是干净参考与带噪目标拼接的视频策略，没有原生联合音频目标。

{\def\LTcaptype{table} % do not increment counter
\small\setlength{\tabcolsep}{3pt}\renewcommand{\arraystretch}{1.18}
\begin{longtable}[]{@{}
  >{\raggedright\arraybackslash}p{(\linewidth - 2\tabcolsep) * \real{0.2200}}
  >{\raggedright\arraybackslash}p{(\linewidth - 2\tabcolsep) * \real{0.7800}}@{}}
\toprule\noalign{}
\begin{minipage}[b]{\linewidth}\raggedright
工程位置
\end{minipage} & \begin{minipage}[b]{\linewidth}\raggedright
候选改动
\end{minipage} \\
\midrule\noalign{}
\endhead
\bottomrule\noalign{}
\endlastfoot
数据加载与collate & 加载参考列表、绑定、源条件、几何、音频状态；按形状与任务分桶 \\
\codepath{ModelInputs}/前向参数 & 另外传入condition bundle；不要借用文本context字段塞入所有类型 \\
块循环 & 在固定12处调用新增支路，保持原块参数与输出头 \\
参数收集 & 同时收集新支路与选定LoRA，冻结codec和文本组件 \\
checkpoint保存 & 除LoRA外保存所有新增模块、门控、规格版本及优化器状态 \\
推理加载 & 先加载对应基座，再按模块名严格加载新增参数和预算配置 \\
\end{longtable}
}

新增模块还要确实进入优化器和保存列表。若它挂在 Transformer 外，原 trainer 只扫描 Transformer 就可能漏掉；挂在内部，也要检查 PEFT 包装和 LoRA 筛选。现有 checkpoint 主要保存模型或 LoRA 权重，并不因此包含 optimizer、随机状态和样本游标。一个 strategy 类只定义前向任务，还没有自动完成这些连接。

\section{哪些结果应能由设计直接解释}\label{seedance-m21-ux54eaux4e9bux7ed3ux679cux5e94ux80fdux7531ux8bbeux8ba1ux76f4ux63a5ux89e3ux91ca}

检查时可以依次关掉门控、替换参考、改变源 mask，再中断恢复。分别观察是否回到原生结构、是否只改变相应绑定、工作区和交付区是否按定义变化、恢复后模块与样本位置是否一致。这些都是实施时需要验证的关系，当前仍没有模型实测结果。

真实实现阅读入口是固定版本的 \href{https://github.com/Lightricks/LTX-2/blob/bd92a5f408eb8c99f8c05e9fe2fcaaec8b56b638/packages/ltx-core/src/ltx_core/model/transformer/model.py}{块循环}、\href{https://github.com/Lightricks/LTX-2/blob/bd92a5f408eb8c99f8c05e9fe2fcaaec8b56b638/packages/ltx-core/src/ltx_core/model/transformer/transformer.py}{联合块} 和 \href{https://github.com/Lightricks/LTX-2/blob/bd92a5f408eb8c99f8c05e9fe2fcaaec8b56b638/packages/ltx-trainer/src/ltx_trainer/trainer.py}{trainer}。新增支路属于待实现的候选结构，不能直接加载为现成上游插件。
